% §13.4 Closing Part III
%
% Monograph register: Part III close + Parts IV/V preview. Fix: "sums go to direct
% sums" -> the sum is the exception (tag coupling). Uses "cipher maps" (not retired
% "oblivious"). Cross-refs ch:approximate-functions, ch:sketches-estimators,
% ch:approximate-relations, sec:composition-functor, ch:optimal-bernoulli-sets,
% ch:perfect-hash-filters, ch:not-bernoulli. DEFINED: sec:closing-part3.

\section{Closing Part III}
\label{sec:closing-part3}

Three parts are now in place. Part~I built the atom: the noisy bit, a
$\Bernoulli{\BoolType}$ value with a $2 \times 2$ channel and rates
$(\fprate, \fnrate)$. Part~II specialized it one direction, to predicates on a universe, and
worked out the algebra of approximate sets, their composition (chapter~6),
classification measures (chapter~7), and information-theoretic floor (chapter~8).
Part~III specialized it a second direction, to maps and types:
\Cref{ch:approximate-functions} lifted the binary channel to $K \times L$
channels, \Cref{ch:sketches-estimators} read streaming counters and estimators
as such channels, chapter~11 placed randomized deciders on the same rate plane,
\Cref{ch:approximate-relations} turned to predicates on pairs, and this chapter
organized the four by type.

Each part widens the type structure. Part~I fixes the codomain to $\BoolType$ and
studies the atom; Part~II keeps the codomain $\BoolType$ and studies functions
$T \to \BoolType$; Part~III releases both domain and codomain to range over the
algebraic types built from $\BoolType$ and the four constructors. The two rates
of Part~I become the channel matrix of Part~III. Under independent noise,
products go to Kronecker products, lists to per-position lifts, functions to
per-input lifts; the sum is the exception, where the tag error couples the
components, the same tag error that is the membership error of an approximate
set. Part~III also recorded two limits that the rest of the book keeps in view:
rates do not propagate as a single pair through computations, and stored stages
or shared inputs break the independence that factorization needs.

Three parts remain. Part~IV turns to constructions near the chapter-8 floor, a
bound that is classical and is met to within lower-order terms by known
retrieval-based filters (chapter~14's notes give the lineage).
\Cref{ch:optimal-bernoulli-sets} analyzes the Bernoulli Hash Function, a
seed-search construction that addresses a fingerprint table through a perfect
hash, and \Cref{ch:perfect-hash-filters} the perfect hash filters built from
minimal perfect hashing. Part~V draws the boundary and turns the algebra to
cryptographic use: \Cref{ch:not-bernoulli} catalogues where the
independent-error model fails (correlated errors, drifting rates, adversarial
queries), and chapters~17 and~18 read the algebra in reverse, the salt of a
Bernoulli hash function as a key and the acceptance predicate as a cipher map.
\Cref{ch:operator-representation} then closes the book by re-reading the
framework in operator form, elements as basis vectors and noise as a stochastic
matrix, where the Kronecker products of \Cref{sec:composition-functor} become
the tensor factorization of an independent set's noise.
